\hmsubsection{System Overview}{FS}{}
This chapter describes how the architecture from
Chapter \ref{cha:system-design} is built in real code and real
hardware. The split is the same as in the architecture chapter: the
Raspberry Pi 5 runs the Python application, the OpenRB-150 runs the
firmware, and the two talk over USB serial.
The Python app is structured as the five layers from
Section \ref{sec:software-architecture}. Frames come in at 20--25 Hz
from the OAK-D S2 camera. The perception layer produces tracked ball detections.
The planning layer turns the chosen target into goal positions for the
five Dynamixel servos. The communication layer sends those goal
positions to the OpenRB-150 and reads back status. The supervisor
runs the state machine.
The firmware on the OpenRB-150 is one C++ program built around the
Arduino \texttt{setup()} and \texttt{loop()} pattern. At
\texttt{setup()} the firmware opens the USB serial port at 115\,200~bps,
opens the Dynamixel bus at 115\,200~bps, and pings each motor. In
\texttt{loop()} it reads commands from the host one line at a time,
sends them to the motors, and reports status back.
The rest of this chapter goes through each part of the implementation
in turn:
\begin{itemize}
 \item Image processing (Section \ref{sec:image-processing}).
 \item Object detection (Section \ref{sec:object-detection}).
 \item Hardware integration (Section \ref{sec:hardware-integration}).
 \item Communication between systems (Section \ref{sec:comm}).
 \item Actuator control (Section \ref{sec:actuator-control}).
 \item Inverse kinematics (Section \ref{sec:ik-solver}).
 \item Calibration subsystem (Section \ref{sec:calibration}).
\end{itemize}
